% ============================================================
% 60_boxes.tex — tcolorbox-Stile, Balken, Boxen, Formelnoten, Listen
% ============================================================

% ------------------------------------------------------------
% Balken-Bindung: ein Titelbalken steht nie allein am Spaltenfuss
% ------------------------------------------------------------
% Ein Balken ohne den Anfang seines Inhalts ist der auffälligste Layoutfehler
% einer ZSF: Man sieht die Überschrift, aber nicht, wozu sie gehört.
%
% Die frühere Platzreserve allein hat das nicht verhindert, sondern verursacht.
% Sie besteht aus \penalty0 … \penalty-100 und legt damit einen BILLIGEN
% Umbruchpunkt an genau die Stelle, an der ein Balken gerade fertig gesetzt
% wurde. Die zugehörige Bremse — Stretch-Glue, die einen Umbruch bei viel
% Restplatz unattraktiv macht — trägt unter \raggedcolumns nicht mehr, weil
% multicols den Spaltenboden ohnehin frei auffüllt. Übrig blieb die Einladung
% zum Umbruch. Genau dort brach die Spalte.
%
% Deshalb zwei getrennte Aufgaben statt einer:
%   \ZSFReserveSpace  — „lieber vorher umbrechen, wenn es eng wird" (Bremse)
%   \ZSFTitleBarLock  — „hier NICHT umbrechen, es kam noch kein Inhalt" (Sperre)
% Die Sperre ist die verlässliche; sie hängt an keiner Höhenschätzung.
\newif\ifzsfTitleBarPending
\newcommand{\ZSFTitleBarLockOn}{\global\zsfTitleBarPendingtrue}
\newcommand{\ZSFTitleBarLockOff}{\global\zsfTitleBarPendingfalse}

% Jeder Umbruchpunkt, den das Balken-System selbst setzt, läuft hierüber:
% solange noch kein Inhalt folgte, wird daraus ein Verbot.
\newcommand{\ZSFBarPenalty}[1]{%
  \ifzsfTitleBarPending\penalty10000\else\penalty#1\fi
}

% Platzreserve vor einem Titelbalken: Ist weniger als #1 frei, bricht die
% Spalte davor um. Genau eine Implementierung — die drei Balkenebenen
% unterscheiden sich nur im Schwellwert (rules/30_spacing).
% Direkt hinter einem Balken entfällt sie zugunsten der Sperre.
\newcommand{\ZSFReserveSpace}[1]{%
  \par
  \ifzsfTitleBarPending
    \penalty10000
  \else
    \penalty0\vspace{0pt plus #1}\penalty-100\vspace{0pt plus -#1}%
  \fi
  \relax
}
\newcommand{\BreakIfNotEnoughSpace}{\ZSFReserveSpace{\ZSFbreakThreshold}}
\newcommand{\BreakIfNotEnoughSpaceChapter}{\ZSFReserveSpace{\ZSFbreakThresholdChapter}}
\newcommand{\BreakIfNotEnoughSpaceSubsubsection}{\ZSFReserveSpace{\ZSFbreakThresholdSubsubsection}}

% Einheitlicher Top-Spacing-Kontrakt für Titelbalken.
% Reihenfolge ist kritisch: \addvspace MUSS vor \BreakIfNotEnoughSpace stehen,
% damit es per Max-Semantik mit dem after-skip der Vorbox kollabiert.
% tcolorbox-Definitionen, die diesen Helper rufen, MUESSEN before skip=0pt
% setzen, damit das Spacing nicht doppelt eingerechnet wird.
\newcommand{\ZSFTitleBarTopSpacing}[1]{%
  \par
  \addvspace{#1}%
  \BreakIfNotEnoughSpace
}

% ------------------------------------------------------------
% Titelbalken — gemeinsame Basis
% ------------------------------------------------------------
% Alle vier Balken (Kapitel, Dokumenttitel, Abschnitt, Unterabschnitt) teilen
% denselben Rahmenlos-Kontrakt. Sie setzen nur noch, worin sie sich wirklich
% unterscheiden: Farbe, vertikale Polsterung, Skips und Schrift.
% Der Balken-Ton ist immer Fläche = Rahmen, deshalb ein Schlüssel dafür.
%
% Nach einem Balken kommt zuerst das Umbruchverbot, dann erst der Abstand.
% Die Reihenfolge ist der eigentliche Fix: tcolorbox' "after skip" setzt die
% Glue direkt hinter die Box, und Glue nach einer Box IST ein gültiger
% Umbruchpunkt. Ein \penalty10000, das erst danach kommt, greift zu spät.
% Steht die Sperre davor, ist die Glue durch ein Penalty abgeschirmt und
% damit kein Umbruchpunkt mehr. Deshalb ersetzen die Balken "after skip"
% durch diesen Hook — "after" ersetzt die Default-Aktion von tcolorbox.
\newcommand{\ZSFTitleBarAfter}[1]{%
  \par\penalty10000\vskip#1\relax\ZSFTitleBarLockOn
}

\tcbset{
  zsfbarcolor/.style={colback=#1, colframe=#1},
  zsfbar/.style={
    enhanced jigsaw,
    breakable=false,
    boxrule=0pt,
    boxsep=0pt,
    left=\ZSFboxPadX,
    right=\ZSFboxPadX,
    before skip=0pt,
  },
}

\newtcolorbox{chapterbar}[1][]{
  zsfbar,
  zsfbarcolor=\ChapterBarColor,
  top=\ZSFbarPadY,
  bottom=\ZSFbarPadY,
  before skip=\ZSFspaceL,
  after={\ZSFTitleBarAfter{\ZSFchapterBarAfterSkip}},
  fontupper=\ZSFfontChapter,
  #1
}
\newcommand{\ChapterBar}[1]{
  \BreakIfNotEnoughSpaceChapter
  \begin{chapterbar}
    \textcolor{\ChapterBarTextColor}{#1}
  \end{chapterbar}
}

% ------------------------------------------------------------
% Dokument-Titel-Masthead — einmalig vor dem Front-Chapter
% ------------------------------------------------------------
\newtcolorbox{zsftitleheader}[1][]{
  zsfbar,
  zsfbarcolor=ZSFDocumentTitleBarColor,
  left=\ZSFboxPadXBar,
  right=\ZSFboxPadXBar,
  top=\ZSFbarPadY,
  bottom=\ZSFbarPadY,
  after={\ZSFTitleBarAfter{\ZSFspaceXS}},
  halign=center,
  #1
}
\newcommand{\ZSFTitleHeader}[2]{%
  \ZSFBarPenalty{-500}%
  \begin{zsftitleheader}
    \sffamily
    {\color{ZSFDocumentTitleTextColor}\ZSFfontTitleHeader #1\par}%
    \vspace{\ZSFspaceXS}%
    {\color{ZSFDocumentTitleBylineColor}\ZSFfontNote von~#2\par}%
  \end{zsftitleheader}
}

% ------------------------------------------------------------
% Subsection-Bar (+ unnummerierte Star-Variante)
% ------------------------------------------------------------
\newtcolorbox{subsectionbar}[1][]{
  zsfbar,
  zsfbarcolor=\SubsectionBarColor,
  top=\ZSFbarPadY,
  bottom=\ZSFbarPadY,
  after={\ZSFTitleBarAfter{\ZSFsubsectionAfterGap}},
  fontupper=\ZSFfontSection,
  #1
}
\newtcolorbox{subsubsectionbar}[1][]{
  zsfbar,
  zsfbarcolor=\SubsubsectionBarColor,
  top=\ZSFbarPadYTight,
  bottom=\ZSFbarPadYTight,
  after={\ZSFTitleBarAfter{\ZSFsubsectionAfterGap}},
  fontupper=\ZSFfontSubsubsection,
  #1
}
\makeatletter
% Pro Ebene EIN Rumpf. Die Varianten liefern nur noch das Nummernpräfix
% (leer bei der Sternvariante) — vorher stand derselbe Rumpf viermal da.
%   #1 = Label (leer erlaubt)   #2 = Nummernpräfix   #3 = Titel
\newcommand{\ZSF@subsectionBarBody}[3]{%
  \ZSFTitleBarTopSpacing{\ZSFsubsectionBarBeforeSkip}%
  \ZSFBarPenalty{-500}%
  \ZSFCollapseAfterChapterBar
  \if\relax\detokenize{#1}\relax\else\label{#1}\fi
  \begin{subsectionbar}
    \textcolor{\SubsectionBarTextColor}{#2#3}
  \end{subsectionbar}
  \ZSFMarkAfterSubsectionBar
}
\newcommand{\ZSF@subsubsectionBarBody}[3]{%
  \BreakIfNotEnoughSpaceSubsubsection
  \par\addvspace{\ZSFsubsubsectionBarBeforeSkip}%
  \ZSFBarPenalty{-300}%
  \if\relax\detokenize{#1}\relax\else\label{#1}\fi
  \begin{subsubsectionbar}
    \textcolor{\SubsubsectionTitleColor}{#2#3}
  \end{subsubsectionbar}
  \ZSFMarkAfterSubsectionBar
}

\newcommand{\SubsectionBar}{\@ifstar{\SubsectionBarStar}{\SubsectionBarNumbered}}
\newcommand{\SubsectionBarNumbered}[2][]{%
  \refstepcounter{subsection}\setcounter{subsubsection}{0}%
  \ZSF@subsectionBarBody{#1}{\thesubsection. }{#2}}
\newcommand{\SubsectionBarStar}[2][]{\ZSF@subsectionBarBody{#1}{}{#2}}

\newcommand{\SubsubsectionBar}{\@ifstar{\SubsubsectionBarStar}{\SubsubsectionBarNumbered}}
\newcommand{\SubsubsectionBarNumbered}[2][]{%
  \refstepcounter{subsubsection}%
  \ZSF@subsubsectionBarBody{#1}{\thesubsubsection. }{#2}}
\newcommand{\SubsubsectionBarStar}[2][]{\ZSF@subsubsectionBarBody{#1}{}{#2}}

% Variante auf neuer Spalte: nur ein \columnbreak davor (ersetzt rohes
% \columnbreak im Kapitel). Kein eigener Rumpf.
\newcommand{\SubsectionBarOnNewColumn}{\@ifstar{\SubsectionBarOnNewColumnStar}{\SubsectionBarOnNewColumnNumbered}}
\newcommand{\SubsectionBarOnNewColumnNumbered}[2][]{\columnbreak\SubsectionBarNumbered[#1]{#2}}
\newcommand{\SubsectionBarOnNewColumnStar}[2][]{\columnbreak\SubsectionBarStar[#1]{#2}}
\newcommand{\SubsubsectionBarOnNewColumn}{\@ifstar{\SubsubsectionBarOnNewColumnStar}{\SubsubsectionBarOnNewColumnNumbered}}
\newcommand{\SubsubsectionBarOnNewColumnNumbered}[2][]{\columnbreak\SubsubsectionBarNumbered[#1]{#2}}
\newcommand{\SubsubsectionBarOnNewColumnStar}[2][]{\columnbreak\SubsubsectionBarStar[#1]{#2}}
\makeatother

% ------------------------------------------------------------
% Gemeinsame tcb-Stile
% ------------------------------------------------------------
\tcbset{zsftitlebar/.style={
  fonttitle=\ZSFfontBoxTitle\strut,
  halign title=center,
  title style={minimum height=\ZSFtitleBarHeight},
  lefttitle=0pt,
  righttitle=0pt,
  toptitle=\ZSFspaceXS,
  bottomtitle=\ZSFspaceXS,
  before title={},
  after title={},
}}

% Rechtsbuendiges Meta-/Kategorie-Tag im Titelbalken einer Titelbox.
% Verwendung im Titel-Argument der Box, z.B.:
%   \begin{defbox}[Quicksort\ZSFTitleTag{O(n log n)}] ... \end{defbox}
%   \begin{codebox}[Algorithmus\ZSFTitleTag{Python}]  ... \end{codebox}
% Schiebt das Tag per \hfill an den rechten Titelrand. Schrift: \ZSFfontTitleTag.
\newcommand{\ZSFTitleTag}[1]{\hfill{\ZSFfontTitleTag #1}}

% Zum "before"-Schlüssel in allen Box-Stilen unten: das führende \par ist Pflicht.
% "before" ERSETZT die Default-Aktion von tcolorbox (die selbst ein \par enthält) —
% es hängt nicht an. Ohne \par wird eine Box, vor der roher Fliesstext ohne
% Leerzeile steht, inline an den offenen Absatz gehängt: sie bekommt keine
% Spaltenbreite und ragt sichtbar in die Nachbarspalte. Im vertikalen Modus ist
% \par ein No-op, der Zusatz ist also folgenlos, wo ohnehin alles korrekt war.
% Gemeinsamer Grundvertrag aller Inhaltsboxen. Was hier steht, gilt fuer jede
% Box — die einzelnen Stile setzen nur noch ihre Unterschiede. Besonders das
% before-Hook ist Pflicht (Begruendung im Kommentar oben) und darf deshalb
% nicht pro Box wiederholt werden, wo man es vergessen kann.
\tcbset{zsfbox/.style={
  enhanced jigsaw,
  before skip=\ZSFBoxBeforeSkip,
  after skip=\ZSFBoxAfterSkip,
  before={\par\ZSFConsumeAfterSubsectionBar\ZSFTitleBarLockOff},
}}

\tcbset{zsfpanelbox/.style={
  zsfbox,
  colback=white,
  colbacktitle=\BoxTitleColor,
  colframe=\BoxTitleColor,
  coltitle=\BoxTitleTextColor,
  boxrule=\ZSFboxRuleWidth,
  titlerule=0pt,
  boxsep=0pt,
  arc=\ZSFboxArc, outer arc=\ZSFboxArc,
  zsftitlebar,
}}

% Single-Point-of-Truth für Titelboxen (defbox/tablebox/figbox/warnbox).
% left/right/top/bottom HIER setzen, damit \linewidth in jeder Titelbox
% identisch ist (sonst variierende Tabellen-Breiten je Box-Typ).
\tcbset{zsftitlebox/.style={
  zsfbox,
  breakable=false,
  colback=\BoxBackColor,
  colbacktitle=\BoxTitleColor,
  colframe=\BoxFrameColor,
  coltitle=\BoxTitleTextColor,
  boxrule=\ZSFboxRuleWidth,
  titlerule=0pt,
  boxsep=0pt,
  arc=\ZSFboxArc,
  outer arc=\ZSFboxArc,
  left=\ZSFboxPadX,
  right=\ZSFboxPadX,
  top=\ZSFboxPadY,
  bottom=\ZSFboxPadY,
  zsftitlebar,
  before upper={\parskip=\ZSFspaceS\relax\ZSFReadableAuto},
}}

% tablebox: zsftitlebox-Rahmen/Ecken, aber ohne Innen-Padding um die Tabelle.
\tcbset{zsfTableboxShell/.style={
  zsftitlebox,
  left=0pt,
  right=0pt,
  top=0pt,
  bottom=0pt,
  colback=white,
  clip upper,
  before upper={\ZSFReadableAuto},
}}

% ------------------------------------------------------------
% Titelboxen
% ------------------------------------------------------------
\newtcolorbox{defbox}[1][]{zsftitlebox, title={#1}}
\newtcolorbox{tablebox}[1][]{zsfTableboxShell, title={#1}}
\newtcolorbox{figbox}[1][]{zsftitlebox, unbreakable, title={#1}}

% Warn-Box (rot) — für Hinweise und Stolperfallen
\newtcolorbox{warnbox}[1][Achtung]{
  zsftitlebox,
  colback=WarnboxBack,
  colbacktitle=ETHRot,
  coltitle=white,
  title=#1,
}

% ------------------------------------------------------------
% Formula-Box
% ------------------------------------------------------------
% Das optionale Argument bleibt eine tcolorbox-Optionsliste. Dadurch bleiben
% bestehende Aufrufe wie \begin{formulabox}[breakable] kompatibel; ein Titel
% wird explizit als \begin{formulabox}[title={Titel}] gesetzt.
\tcbset{zsfformulastyle/.style={
  zsfbox,
  colback=\FormulaboxBackColor,
  colbacktitle=\FormulaboxBackColor,
  colframe=black,
  coltitle=TextVorBoxColor,
  boxrule=\ZSFboxRuleWidth,
  titlerule=0pt,
  fonttitle=\ZSFfontBoxTitle,
  halign title=center,
  lefttitle=\ZSFboxPadX,
  righttitle=\ZSFboxPadX,
  toptitle=\ZSFspaceXS,
  bottomtitle=0pt,
  left=\ZSFboxPadX,
  right=\ZSFboxPadX,
  top=\ZSFboxPadY,
  bottom=\ZSFboxPadY,
  before upper={\parskip=\ZSFspaceS\centering},
  after upper={\par},
  middle=\ZSFspaceSep,
}}
\newtcolorbox{formulabox}[1][]{zsfformulastyle,#1}

% Formel-Note (klein, kursiv, dezent) — unter oder über Formeln
\newcommand{\formulanote}[1]{%
  \ZSFInterlude{\ZSFspaceS}%
  \noindent{{\ZSFfontNote #1}}%
  \ZSFInterlude{\ZSFspaceS}%
}
% \ZSFformulaline{<mathe>}{<Anmerkung>} — Formelzeile mit rechtsbündiger
% Kurz-Anmerkung auf derselben Höhe (Gültigkeitsbereich, Einheit, Fallname).
% \formulanote setzt die Anmerkung darunter; diese Form hält sie auf Zeile.
\newcommand{\ZSFformulaline}[2]{%
  \par\noindent$\displaystyle #1$\hfill{\ZSFfontNote #2}\par
}

% Duenne, graue Trennlinie zwischen Formeln in einer formulabox
\newcommand{\formulasep}{%
  \ZSFInterlude{\ZSFspaceSep}%
  \noindent\textcolor{FormulaSeparatorColor}{\rule{\linewidth}{\ZSFboxRuleWidth}}%
  \ZSFInterlude{\ZSFspaceSep}%
  \penalty-50%
}

% Zentriertes Item-Heading mit Trennstrich — für gleichwertige Fälle in Boxen
\newcommand{\ZSFItemHeading}[1]{%
  \ZSFInterlude{\ZSFspaceM}%
  {\centering\textbf{#1}\par}%
  \ZSFInterlude{\ZSFspaceS}%
  \noindent\rule{\linewidth}{\ZSFboxRuleWidth}%
  \ZSFInterlude{\ZSFspaceS}%
}

% ------------------------------------------------------------
% Runintext + Grafik-Helpers
% ------------------------------------------------------------
\newenvironment{runintext}{%
  \ZSFConsumeAfterChapterBar\ZSFTitleBarLockOff
  \par\addvspace{\ZSFspaceS}\noindent\ZSFReadableAuto\ignorespaces
}{\par}

% Dezente Bildunterschrift (leerer Text -> nichts)
\newcommand{\ZSFfigcaption}[1]{%
  \if\relax\detokenize{#1}\relax\else
    \par\vspace{\ZSFspaceXS}%
    {\ZSFfontNote\color{MutedTextColor}#1\par}%
  \fi
}

% \ZSFfig[breite]{pfad}{Titel}{Caption} — figbox mit Bild + optionaler Caption.
% Ohne [breite]: min(Spaltenbreite, ZSFfigMaxHeight). Mit einem Anteil von
% \linewidth (z.B. [0.7]): exakte lokale Breite ohne Höhen-Cap.
\newcommand{\ZSFfig}[4][]{%
  \begin{figbox}[{#3}]%
    \centering
    \if\relax\detokenize{#1}\relax
      \includegraphics[width=\linewidth,height=\ZSFfigMaxHeight,keepaspectratio]{#2}%
    \else
      \includegraphics[width=\dimexpr#1\linewidth\relax,keepaspectratio]{#2}%
    \fi
    \ZSFfigcaption{#4}%
  \end{figbox}%
}

% \ZSFfigrow{pfad1,pfad2,...}{Titel}{Caption} — mehrere Abbildungen
% gleicher Breite nebeneinander in EINER figbox. Die Breite ergibt sich
% aus der Anzahl Pfade; der Rest verteilt sich als Zwischenraum.
% Ersetzt handgebaute \includegraphics-Reihen mit \hfill.
\makeatletter
\newcounter{ZSF@figrowNum}
\newlength{\ZSF@figrowWidth}
\newcommand{\ZSFfigrow}[3]{%
  \begin{figbox}[{#2}]%
    \setcounter{ZSF@figrowNum}{0}%
    \@for\ZSF@figrowItem:=#1\do{\stepcounter{ZSF@figrowNum}}%
    \setlength{\ZSF@figrowWidth}{%
      \dimexpr0.97\linewidth/\value{ZSF@figrowNum}\relax}%
    \noindent
    \@for\ZSF@figrowItem:=#1\do{%
      \begin{minipage}[c]{\ZSF@figrowWidth}%
        \centering
        \includegraphics[width=\linewidth,height=\ZSFfigMaxHeight,%
          keepaspectratio]{\ZSF@figrowItem}%
      \end{minipage}%
      \hfill
    }%
    \unskip
    \par
    \ZSFfigcaption{#3}%
  \end{figbox}%
}
\makeatother

% \ZSFfigside[Bildanteil]{pfad}{Titel}{Inhalt rechts} — Bild links,
% Inhalt rechts. Der Anteil wird erst am Einsatzort gegen die dortige
% \linewidth ausgewertet; Standard sind 30%.
\newcommand{\ZSFfigside}[4][0.30]{%
  \begin{figbox}[{#3}]%
    \noindent
    \begin{minipage}[c]{#1\linewidth}%
      \centering
      \includegraphics[width=\linewidth,height=\ZSFfigMaxHeight,keepaspectratio]{#2}%
    \end{minipage}%
    \hfill
    \begin{minipage}[c]{\dimexpr0.96\linewidth-#1\linewidth\relax}%
      \raggedright
      #4%
    \end{minipage}%
    \par
  \end{figbox}%
}

% Side-by-side-Layout (Bild + Text in gleicher Linie). Das zweite optionale
% Argument nimmt tcolorbox-Schlüssel auf, z.B. [sidebyside align=center].
\NewDocumentEnvironment{splitbox}{O{0.3} O{}}{%
  \begin{tcolorbox}[
    blankest,
    sidebyside,
    sidebyside align=top,
    sidebyside gap=\ZSFspaceM,
    lefthand width=#1\linewidth,
    lower separated=false,
    before={\par\ZSFConsumeAfterSubsectionBar\ZSFTitleBarLockOff},
    before skip=\ZSFspaceS,
    after skip=\ZSFspaceS,
    before upper={\ZSFReadableAuto},
    before lower={\ZSFReadableAuto},
    boxsep=0pt, left=0pt, right=0pt, top=0pt, bottom=0pt,
    #2
  ]%
}{%
  \end{tcolorbox}%
}

% ------------------------------------------------------------
% Aussagen / Verfahren / Listen
% ------------------------------------------------------------
% statementbox: dezent, linker Farbbalken in Kapitelfarbe.
\tcbset{zsfstatementstyle/.style={
  zsfbox,
  breakable=false,
  colback=white,
  colframe=white,
  frame hidden,
  borderline west={\ZSFstatementBarWidth}{0pt}{\StatementBarColor},
  boxrule=0pt,
  arc=0pt, outer arc=0pt,
  left=\ZSFboxPadXBar,
  right=\ZSFboxPadX,
  top=\ZSFboxPadYTight,
  bottom=\ZSFboxPadYTight,
  before upper={\ZSFReadableAuto},
}}

\NewDocumentEnvironment{statementbox}{O{}}
  {\begin{tcolorbox}[zsfstatementstyle]%
   \if\relax\detokenize{#1}\relax\else
     {\bfseries\color{\StatementTitleColor}#1}\par\vspace{\ZSFspaceXS}%
   \fi\ignorespaces}
  {\end{tcolorbox}}

% procedure: nummerierte Schritt-Liste in einer statementbox.
\NewDocumentEnvironment{procedure}{O{}}
  {\begin{statementbox}[#1]%
   \begin{enumerate}[leftmargin=\ZSFprocLeftMargin,labelsep=\ZSFlistLabelSep,itemsep=\ZSFprocItemSep,topsep=0pt,parsep=0pt,partopsep=0pt]}
  {\end{enumerate}\end{statementbox}}
\newcommand{\ProcStep}[2]{\item {\bfseries #1}\\[\ZSFprocStepGap]#2}

% factlist: Liste von Definitions-/Eigenschaftsfakten, Einträge via \ZSFFact.
% Ohne Titel rahmenlos, mit Titel in einer statementbox — dieselbe Liste,
% eine Option statt zweier Environments. Das Aufzählungszeichen bleibt in
% beiden Fällen identisch; früher unterschieden sie sich hier grundlos.
\makeatletter
\newif\ifZSF@factlistBoxed
\newcommand{\ZSF@factlistOpen}[1]{%
  \begin{itemize}[
    leftmargin=\ZSFlistLeftMargin,
    labelsep=\ZSFlistLabelSep,
    label={\textcolor{\StatementBarColor}{\textbullet}},
    itemsep=\ZSFprocItemSep,
    topsep=#1,
    parsep=0pt,
    partopsep=0pt
  ]}
\NewDocumentEnvironment{factlist}{o}
  {\IfNoValueTF{#1}%
     {\ZSF@factlistBoxedfalse\ZSF@factlistOpen{\ZSFspaceXS}}%
     {\ZSF@factlistBoxedtrue\begin{statementbox}[#1]\ZSF@factlistOpen{0pt}}}
  {\end{itemize}\ifZSF@factlistBoxed\end{statementbox}\fi} % chktex 9
\makeatother
\newcommand{\ZSFFact}[2]{\item \textbf{#1}\,---\,#2}

% ------------------------------------------------------------
% Goal-System: Bedingung -> Umformung -> Ziel (Inline + Card)
% ------------------------------------------------------------
\newtcolorbox{goalbox}[1][]{
  zsfpanelbox,
  breakable=false,
  colback=white,
  left=\ZSFspaceS,
  right=\ZSFspaceS,
  top=\ZSFtextMinPad,
  bottom=\ZSFtextMinPad,
  title={#1},
  before upper={\parskip=\ZSFspaceS\centering},
  after upper={\par},
}

\newtcbox{\GoalConditionBox}{%
  on line, enhanced,
  boxrule=\ZSFgoalConditionRuleWidth,
  colback=\GoalConditionBackColor,
  colframe=\GoalConditionFrameColor,
  borderline={\ZSFgoalConditionBorderlineWidth}{0pt}{black},
  arc=\ZSFboxArc,
  left=\ZSFspaceXS, right=\ZSFspaceXS, top=\ZSFspaceXS, bottom=\ZSFspaceXS,
  nobeforeafter,
  tcbox raise base,
}
\newtcbox{\GoalTargetBox}{%
  on line, enhanced,
  boxrule=\ZSFboxRuleWidth,
  colback=\GoalTargetBackColor,
  colframe=\GoalTargetFrameColor,
  arc=\ZSFboxArc,
  left=\ZSFspaceXS, right=\ZSFspaceXS, top=\ZSFspaceXS, bottom=\ZSFspaceXS,
  nobeforeafter,
  tcbox raise base,
}

\newcommand{\GoalCondition}[1]{\par\centering\GoalConditionBox{\ZSFfontNote #1}\par}
\newcommand{\GoalStep}[1]{\par\centering$\displaystyle #1$\par}
\newcommand{\GoalTarget}[1]{\par\centering\GoalTargetBox{$\displaystyle #1$}\par}
\newcommand{\ZSFDerivationCase}[3]{%
  \GoalCondition{#1}%
  \GoalStep{#2 =}%
  \GoalTarget{#3}%
}
\newcommand{\ZSFDerivationInlineCase}[3]{%
  \GoalCondition{#1}%
  \GoalStep{#2 = \GoalTargetBox{$\displaystyle #3$}}%
}
% Inline-Danger-Tag (rotes Pill-Label für Stolperfallen im Fliesstext)
\newtcbox{\ZSFdanger}{
  colback=DangerTagColor, colframe=DangerTagColor, coltext=white,
  on line, boxsep=\ZSFspaceXS,
  left=\ZSFspaceS, right=\ZSFspaceS, top=\ZSFspaceXS, bottom=\ZSFspaceXS,
  boxrule=0pt, arc=\ZSFboxArc,
  fontupper=\sffamily\bfseries
}

% ------------------------------------------------------------
% Valuegrid: kompakte (tabularray-)Wertegrids in einer Panel-Box
% ------------------------------------------------------------
% Interner Panel-Rahmen des Valuegrids — kein öffentlicher Baustein.
% Kapitel greifen ausschliesslich über \begin{valuegrid}{n}[Titel] zu.
\tcbset{zsfvaluegridshell/.style={
  zsfpanelbox,
  colframe=NeutralBoxFrameColor,
  breakable=false,
  boxsep=0pt, left=0pt, right=0pt, top=\ZSFtextMinPad, bottom=\ZSFtextMinPad,
  zsftitlebar,
  boxrule=\ZSFboxRuleWidth
}}
\newcommand{\ZSFvaluegridCommon}{
  colsep=\ZSFtableEdgePad,
  rowsep=\ZSFvaluegridRowSep,
  hlines,
  vlines,
  rows={valign=m}
}
% \begin{valuegrid}{<n_datacols>}[<title>] — Spec wird automatisch generiert.
\NewDocumentEnvironment{valuegrid}{m O{}}
  {\begin{tcolorbox}[zsfvaluegridshell, title={#2}]%
   \begin{tblr}{colspec={l *{#1}{c}},\ZSFvaluegridCommon}}
  {\end{tblr}\end{tcolorbox}}

% ------------------------------------------------------------
% Pack-Modus für Anhang-artige Kapitel
% ------------------------------------------------------------
% Ersetzen rohes \tcbset und \renewcommand im Kapitel — beides ist
% Systemumbau und bricht den Linter (rules/02_ai_mandate). Am Anfang des
% betreffenden Kapitels aufrufen; die Wirkung endet mit dessen Gruppe.
%
% Wichtig: Ein globales \tcbset{breakable} im Kapitel hätte ohnehin nicht
% gewirkt — die Basisstile setzen breakable=false in ihrer eigenen
% Optionsliste und gewinnen damit gegen den Default. Deshalb hängen diese
% Schalter die Option an die Stile selbst an.
%
% Bewusst NICHT breakable: figbox, valuegrid, goalbox und die Balken.
% Ein Umbruch mitten in Bild, Wertegrid oder Zielkette ergibt keinen Sinn.
\newcommand{\ZSFBoxesBreakableOn}{%
  \tcbset{zsftitlebox/.append style={breakable},
          zsfstatementstyle/.append style={breakable}}%
}
\newcommand{\ZSFBoxesBreakableOff}{%
  \tcbset{zsftitlebox/.append style={breakable=false},
          zsfstatementstyle/.append style={breakable=false}}%
}
% Hebt die Platzreserven vor Boxen und Überschriften auf, damit ein
% Anhang dicht packt statt früh umzubrechen.
\newcommand{\ZSFBreakReservesOff}{%
  \renewcommand{\BreakIfNotEnoughSpace}{}%
  \renewcommand{\BreakIfNotEnoughSpaceChapter}{}%
  \renewcommand{\BreakIfNotEnoughSpaceSubsubsection}{}%
}
